% GENERATED FILE. Edit canonical lessons, then run npm run generate:books.
% canonical-source-hash: fnv1a64:9416b966dae9d021
% canonical-lessons: GU-C354-adhikari, GU-C354-mantri, GU-C354-manjuri, GU-C354-parvangi, GU-C354-nimvu, GU-R354-review-347-354

\chapter{Offices and Officials}
\label{ch:gu-offices-and-officials}
\hlchaptermodality{354}

\begin{chapteropening}
\textbf{By the end of this chapter:} \emph{I can deal with an office and its officials.}

Complete the last lesson of chapter 354: I can deal with an office and its officials.
\end{chapteropening}

\section[\texorpdfstring{\gu{અધિકારી} (adhikārī)}{અધિકારી (adhikārī)}]{\gu{અધિકારી} (adhikārī) — an official}

\label{lesson:GU-C354-adhikari}



\begin{quote}
Before the new one: say the Gujarati for a candidate, then the Gujarati for to be elected.
\end{quote}

\subsection*{You'll want to know: \gu{અધિકારી}}
\textbf{\gu{અધિકારી}} — \emph{adhikārī} — \textquotedblleft{}an official\textquotedblright{}.

\begin{quote}
\gu{કચેરીમાં} \gu{એક} \gu{અધિકારીએ} \gu{મારી} \gu{અરજી} \gu{તપાસી}.
\end{quote}

\begin{quote}
\emph{kacerīmā̃ ek adhikārīe mārī arjī tapāsī.}
\end{quote}

\begin{quote}
An official at the office checked my application.
\end{quote}

A masculine or feminine noun, by the person meant.

\begin{cousinweb}
\textbf{\gu{અધિકારી}} (\emph{adhikārī}) is a tatsama, borrowed from Sanskrit \emph{adhikārin}, \textquotedblleft{}having authority\textquotedblright{} (Monier-Williams). It shares its source with \textbf{\gu{અધિકાર}} (\emph{adhikār}), \textquotedblleft{}authority\textquotedblright{}, but names the person who holds it.
\end{cousinweb}

\begin{grammarlens}[title={What you've built}]
One more word for giving reasons and opinions.
\end{grammarlens}

\subsection*{Your turn}
\begin{itemize}
\raggedright
  \item \emph{Say it:} \emph{adhikārī}
  \item \emph{Say it:} \emph{adhikārī}, once more
  \item \emph{Say it:} \emph{kacerīmā̃ ek adhikārīe mārī arjī tapāsī.} — An official at the office checked my application.
  \item \emph{Recall:} say the Gujarati for a candidate, then the Gujarati for to be elected, then say \emph{adhikārī} again
\end{itemize}

\begin{tcolorbox}[breakable,colback=teal!4,colframe=teal!35!black,title={Before you move on}]
What does \gu{અધિકારી} mean? (\textquotedblleft{}An official\textquotedblright{}.) Say it once more.
\end{tcolorbox}
\section[\texorpdfstring{\gu{મંત્રી} (mantrī)}{મંત્રી (mantrī)}]{\gu{મંત્રી} (mantrī) — a minister}

\label{lesson:GU-C354-mantri}



\begin{quote}
Before the new one: say the Gujarati for to be elected, then the Gujarati for an official.
\end{quote}

\subsection*{You'll want to know: \gu{મંત્રી}}
\textbf{\gu{મંત્રી}} — \emph{mantrī} — \textquotedblleft{}a minister\textquotedblright{}.

\begin{quote}
\gu{મંત્રી} \gu{કાલે} \gu{અમારા} \gu{ગામની} \gu{મુલાકાત} \gu{લેશે}.
\end{quote}

\begin{quote}
\emph{mantrī kāle amārā gāmnī mulākāt leśe.}
\end{quote}

\begin{quote}
The minister will visit our village tomorrow.
\end{quote}

A masculine or feminine noun, by the person meant.

\begin{cousinweb}
\textbf{\gu{મંત્રી}} (\emph{mantrī}) is a tatsama, borrowed from Sanskrit \emph{mantrin}, \textquotedblleft{}wise; a counsellor, a minister\textquotedblright{} (Monier-Williams).
\end{cousinweb}

\begin{grammarlens}[title={What you've built}]
One more word for giving reasons and opinions.
\end{grammarlens}

\subsection*{Your turn}
\begin{itemize}
\raggedright
  \item \emph{Say it:} \emph{mantrī}
  \item \emph{Say it:} \emph{mantrī}, once more
  \item \emph{Say it:} \emph{mantrī kāle amārā gāmnī mulākāt leśe.} — The minister will visit our village tomorrow.
  \item \emph{Recall:} say the Gujarati for to be elected, then the Gujarati for an official, then say \emph{mantrī} again
\end{itemize}

\begin{tcolorbox}[breakable,colback=teal!4,colframe=teal!35!black,title={Before you move on}]
What does \gu{મંત્રી} mean? (\textquotedblleft{}A minister\textquotedblright{}.) Say it once more.
\end{tcolorbox}
\section[\texorpdfstring{\gu{મંજૂરી} (manjūrī)}{મંજૂરી (manjūrī)}]{\gu{મંજૂરી} (manjūrī) — approval}

\label{lesson:GU-C354-manjuri}



\begin{quote}
Before the new one: say the Gujarati for an official, then the Gujarati for a minister.
\end{quote}

\subsection*{You'll want to know: \gu{મંજૂરી}}
\textbf{\gu{મંજૂરી}} — \emph{manjūrī} — \textquotedblleft{}approval\textquotedblright{}.

\begin{quote}
\gu{સરકારે} \gu{આ} \gu{યોજનાને} \gu{મંજૂરી} \gu{આપી}.
\end{quote}

\begin{quote}
\emph{sarkāre ā yojnāne manjūrī āpī.}
\end{quote}

\begin{quote}
The government approved this plan.
\end{quote}

A feminine noun.

\begin{cousinweb}
\textbf{\gu{મંજૂરી}} (\emph{manjūrī}) is usually traced to Persian \emph{manzūrī}, \textquotedblleft{}approval\textquotedblright{}, from Arabic \emph{manẓūr}, \textquotedblleft{}seen, looked upon with favour\textquotedblright{}.
\end{cousinweb}

\begin{grammarlens}[title={What you've built}]
One more word for giving reasons and opinions.
\end{grammarlens}

\subsection*{Your turn}
\begin{itemize}
\raggedright
  \item \emph{Say it:} \emph{manjūrī}
  \item \emph{Say it:} \emph{manjūrī}, once more
  \item \emph{Say it:} \emph{sarkāre ā yojnāne manjūrī āpī.} — The government approved this plan.
  \item \emph{Recall:} say the Gujarati for an official, then the Gujarati for a minister, then say \emph{manjūrī} again
\end{itemize}

\begin{tcolorbox}[breakable,colback=teal!4,colframe=teal!35!black,title={Before you move on}]
What does \gu{મંજૂરી} mean? (\textquotedblleft{}Approval\textquotedblright{}.) Say it once more.
\end{tcolorbox}
\section[\texorpdfstring{\gu{પરવાનગી} (parvāngī)}{પરવાનગી (parvāngī)}]{\gu{પરવાનગી} (parvāngī) — permission}

\label{lesson:GU-C354-parvangi}



\begin{quote}
Before the new one: say the Gujarati for a minister, then the Gujarati for approval.
\end{quote}

\subsection*{You'll want to know: \gu{પરવાનગી}}
\textbf{\gu{પરવાનગી}} — \emph{parvāngī} — \textquotedblleft{}permission\textquotedblright{}.

\begin{quote}
\gu{અંદર} \gu{જવા} \gu{માટે} \gu{ચોકીદારની} \gu{પરવાનગી} \gu{લેવી} \gu{પડશે}.
\end{quote}

\begin{quote}
\emph{andar javā māṭe cokīdārnī parvāngī levī paḍśe.}
\end{quote}

\begin{quote}
You will need the watchman's permission to go in.
\end{quote}

A feminine noun.

\begin{cousinweb}
\textbf{\gu{પરવાનગી}} (\emph{parvāngī}) is usually traced to Persian \emph{parvānagī}, \textquotedblleft{}permission\textquotedblright{}, from \emph{parvāna}, \textquotedblleft{}a written order\textquotedblright{}.
\end{cousinweb}

\begin{grammarlens}[title={What you've built}]
One more word for giving reasons and opinions.
\end{grammarlens}

\subsection*{Your turn}
\begin{itemize}
\raggedright
  \item \emph{Say it:} \emph{parvāngī}
  \item \emph{Say it:} \emph{parvāngī}, once more
  \item \emph{Say it:} \emph{andar javā māṭe cokīdārnī parvāngī levī paḍśe.} — You will need the watchman's permission to go in.
  \item \emph{Recall:} say the Gujarati for a minister, then the Gujarati for approval, then say \emph{parvāngī} again
\end{itemize}

\begin{tcolorbox}[breakable,colback=teal!4,colframe=teal!35!black,title={Before you move on}]
What does \gu{પરવાનગી} mean? (\textquotedblleft{}Permission\textquotedblright{}.) Say it once more.
\end{tcolorbox}
\section[\texorpdfstring{\gu{નીમવું} (nīmvũ)}{નીમવું (nīmvũ)}]{\gu{નીમવું} (nīmvũ) — to appoint}

\label{lesson:GU-C354-nimvu}



\begin{quote}
Before the new one: say the Gujarati for approval, then the Gujarati for permission.
\end{quote}

\subsection*{You'll want to know: \gu{નીમવું}}
\textbf{\gu{નીમવું}} — \emph{nīmvũ} — \textquotedblleft{}to appoint\textquotedblright{}.

\begin{quote}
\gu{સરકાર} \gu{નવા} \gu{અધિકારીને} \gu{નીમશે}.
\end{quote}

\begin{quote}
\emph{sarkār navā adhikārīne nīmśe.}
\end{quote}

\begin{quote}
The government will appoint a new official.
\end{quote}

Past, with \gu{તેણે} (he or she …): \gu{તેણે} \gu{નવો} \gu{ચોકીદાર} \gu{નીમ્યો}, \emph{teṇe navo cokīdār nīmyo}.

\begin{cousinweb}
\textbf{\gu{નીમવું}} (\emph{nīmvũ}) is inherited from Sanskrit \emph{niyamayati}, \textquotedblleft{}restrains\textquotedblright{}; CDIAL lists the Gujarati verb, \textquotedblleft{}to fix, to appoint\textquotedblright{} (CDIAL 7251). \textbf{\gu{નિયમ}} (\emph{niyam}), \textquotedblleft{}a rule\textquotedblright{}, is its learned cousin.
\end{cousinweb}

\begin{grammarlens}[title={What you've built}]
One more word for giving reasons and opinions.
\end{grammarlens}

\subsection*{Your turn}
\begin{itemize}
\raggedright
  \item \emph{Say it:} \emph{nīmvũ}
  \item \emph{Say it:} \emph{nīmvũ}, once more
  \item \emph{Say it:} \emph{sarkār navā adhikārīne nīmśe.} — The government will appoint a new official.
  \item \emph{Recall:} say the Gujarati for approval, then the Gujarati for permission, then say \emph{nīmvũ} again
\end{itemize}

\begin{tcolorbox}[breakable,colback=teal!4,colframe=teal!35!black,title={Before you move on}]
What does \gu{નીમવું} mean? (\textquotedblleft{}To appoint\textquotedblright{}.) Say it once more.
\end{tcolorbox}
\section[(dialogue)]{Looking back: from The Law to Offices and Officials}

\label{lesson:GU-R354-review-347-354}



\begin{quote}
Say the Gujarati for a law and a violation.
\end{quote}

\subsection*{Your turn}
Say these aloud:

\begin{itemize}
\raggedright
  \item a law, a violation, observance, in force, to defy
  \item a judge, a court case, an appeal, a ruling, to hit hard
  \item reasonable, discrimination, equal, neutral, to behave
  \item a contract, a document, free of, seized, to recover a debt
  \item a fact, an illusion, a misunderstanding, true to the facts, to be taken in
  \item a government, a state (of a country), power to rule, the people of a state, to shape
  \item a citizen, the public, voting, a candidate, to be elected
  \item an official, a minister, approval, permission, to appoint
\end{itemize}

\begin{tcolorbox}[breakable,colback=teal!4,colframe=teal!35!black,title={Before you move on}]
A law? (\textbf{\gu{કાયદો}}, \emph{kāydo}.) A violation? (\textbf{\gu{ઉલ્લંઘન}}, \emph{ullanghan}.) Observance? (\textbf{\gu{પાલન}}, \emph{pālan}.) In force? (\textbf{\gu{લાગુ}}, \emph{lāgu}.) To defy? (\textbf{\gu{ઉથાપવું}}, \emph{uthāpvũ}.) A judge? (\textbf{\gu{ન્યાયાધીશ}}, \emph{nyāyādhīś}.) A court case? (\textbf{\gu{કેસ}}, \emph{kes}.) An appeal? (\textbf{\gu{અપીલ}}, \emph{apīl}.) A ruling? (\textbf{\gu{ફેંસલો}}, \emph{phẽslo}.) To hit hard? (\textbf{\gu{ફટકારવું}}, \emph{phaṭkārvũ}.) Reasonable? (\textbf{\gu{વાજબી}}, \emph{vājbī}.) Discrimination? (\textbf{\gu{ભેદભાવ}}, \emph{bhedbhāv}.) Equal? (\textbf{\gu{સમાન}}, \emph{samān}.) Neutral? (\textbf{\gu{તટસ્થ}}, \emph{taṭasth}.) To behave? (\textbf{\gu{વર્તવું}}, \emph{vartvũ}.) A contract? (\textbf{\gu{કરાર}}, \emph{karār}.) A document? (\textbf{\gu{દસ્તાવેજ}}, \emph{dastāvej}.) Free of? (\textbf{\gu{મુક્ત}}, \emph{mukt}.) Seized? (\textbf{\gu{જપ્ત}}, \emph{japt}.) To recover a debt? (\textbf{\gu{વસૂલવું}}, \emph{vasūlvũ}.) A fact? (\textbf{\gu{હકીકત}}, \emph{hakīkat}.) An illusion? (\textbf{\gu{ભ્રમ}}, \emph{bhram}.) A misunderstanding? (\textbf{\gu{ગેરસમજ}}, \emph{gersamaj}.) True to the facts? (\textbf{\gu{યથાર્થ}}, \emph{yathārth}.) To be taken in? (\textbf{\gu{છેતરાવું}}, \emph{chetrāvũ}.) A government? (\textbf{\gu{સરકાર}}, \emph{sarkār}.) A state (of a country)? (\textbf{\gu{રાજ્ય}}, \emph{rājya}.) Power to rule? (\textbf{\gu{સત્તા}}, \emph{sattā}.) The people of a state? (\textbf{\gu{પ્રજા}}, \emph{prajā}.) To shape? (\textbf{\gu{ઘડવું}}, \emph{ghaḍvũ}.) A citizen? (\textbf{\gu{નાગરિક}}, \emph{nāgrik}.) The public? (\textbf{\gu{જનતા}}, \emph{jantā}.) Voting? (\textbf{\gu{મતદાન}}, \emph{matdān}.) A candidate? (\textbf{\gu{ઉમેદવાર}}, \emph{umedvār}.) To be elected? (\textbf{\gu{ચૂંટાવું}}, \emph{cū̃ṭāvũ}.) An official? (\textbf{\gu{અધિકારી}}, \emph{adhikārī}.) A minister? (\textbf{\gu{મંત્રી}}, \emph{mantrī}.) Approval? (\textbf{\gu{મંજૂરી}}, \emph{manjūrī}.) Permission? (\textbf{\gu{પરવાનગી}}, \emph{parvāngī}.) To appoint? (\textbf{\gu{નીમવું}}, \emph{nīmvũ}.)
\end{tcolorbox}
